\documentclass[../../main]{subfiles}

\ifSubfilesClassLoaded{%
    \zexternaldocument*{../1.Category-Theory/exported/aux/main}%
    \zexternaldocument*{1.Additive-categories}%
}{}

\begin{document}
\chapter{Abelian categories}

In this chapter we introduce abelian categories: \hyperref[def:additive-category]{additive categories} where the first isomorphism theorem works.
We also study exactness in these categories and introduce the cohomology functor.

The main references for this chapter are \cite{weibelIntroductionHomologicalAlgebra1994, gelfandMethodsHomologicalAlgebra2003, aluffiAlgebraChapter02009}.

Before we begin, we fix some terminology and notation around the necessary building blocks of the isomorphism theorem.
\begin{definition}
  Let \(f\colon X \to Y\) be a morphism in an \hyperref[def:additive-category]{additive category} \(\category{C}\) that admits \hyperref[def:categorical-kernel]{kernels} and \hyperref[def:categorical-cokernel]{cokernels}.
  Consider the following
  \begin{center}
    \begin{tikzcd}
      \ker f \arrow[r, "h"] & X \arrow[r, "f"] & Y \arrow[r, "k"] & \coker f
    \end{tikzcd}
  \end{center}
  We denote
  \begin{equation*}
    \operatorname{Im} f \coloneq \ker k \qquad \operatorname{Coim} f \coloneq \coker h
  \end{equation*}
  I.e: we have the following diagram
  \begin{center}
    \begin{tikzcd}
      \ker f \arrow[r, "h"] & X \arrow[r, "f"] \arrow[d, "s"'] & Y \arrow[r, "k"] & \coker f \\
      & \operatorname{Coim} f \arrow[ur, dotted, "\widetilde{f}"] \arrow[r, dotted, "u"'] & \operatorname{Im} f \arrow[u]
    \end{tikzcd}
  \end{center}
  Note that, since \(f \circ h = 0\), \(f\) factors uniquely through \(\operatorname{Coim} f\) producing the morphism \(\widetilde{f}\).
  Furthermore since \(k \circ f = k \circ \widetilde{f} = 0\) we have that \(\widetilde{f}\) factors uniquely through \(\operatorname{Im} f\).
  Therefore we have a canonical morphism
  \begin{equation*}
    u\colon \operatorname{Coim} f \to \operatorname{Im} f
  \end{equation*}
\end{definition}

Now we are ready to define our categories of interest.
\begin{definition}[Abelian category]\label{def:abelian-category}
  An \hyperref[def:additive-category]{additive category} \(\category{C}\) is \textbf{abelian} if it verifies
  \begin{enumerate}[label=Ab.\arabic*), ref=Ab.\arabic*]
    \item Any morphism in \(\category{C}\) admits a \hyperref[def:categorical-kernel]{kernel} and a \hyperref[def:categorical-cokernel]{cokernel}.
    \item For any morphism \(f\) in \(\category{C}\) the canonical morphism \(\operatorname{Coim} f \to \operatorname{Im} f\) is an \hyperref[def:categorical-isomorphism]{isomorphism}.
  \end{enumerate}
\end{definition}

\begin{example}
  The \hyperref[def:category]{category} of (abelian) groups and modules over a ring are \hyperref[def:abelian-category]{abelian}.
  However, the \hyperref[def:category]{category} of rings is not.
\end{example}


As a direct consecuence of \autoref{thm:existence-finite-limits-colimits} we obtain that \hyperref[def:categorical-limit]{limits} and \hyperref[def:categorical-colimit]{colimits} exist in these \hyperref[def:category]{categories}
\begin{lemma}\label{lemma:abelian-admits-limits-colimits}
  An \hyperref[def:abelian-category]{abelian category} admits all finite \hyperref[def:categorical-limit]{limits} and \hyperref[def:categorical-colimit]{colimits}.
\end{lemma}

In the category of modules, an easy way to check whether homomorphisms are \hyperref[def:mono-epi-morphism]{monomorphisms} (resp. \hyperref[def:mono-epi-morphism]{epimorphisms}) is to check their \hyperref[def:categorical-kernel]{kernels} (resp. \hyperref[def:categorical-kernel]{cokernels}).
This still works in \hyperref[def:abelian-category]{abelian categories}.
\begin{lemma}\label{lemma:partial-sequence-abelian}
  Let \(f\colon X \to Y\) be a morphism in an \hyperref[def:abelian-category]{abelian category}. Then \(f\) is a \hyperref[def:mono-epi-morphism]{monomorphism} if and only if \(0 \to X\) is its \hyperref[def:categorical-kernel]{kernel}.

  Dually, \(f\) is an \hyperref[def:mono-epi-morphism]{epimorphism} if and only if \(Y \to 0\) is its \hyperref[def:categorical-cokernel]{cokernel}.
\end{lemma}
\begin{proof}
  We will do the \hyperref[def:categorical-kernel]{kernel} one.
  \begin{itemize}
    \item[\(\implies\)] Suppose \(f\) is a \hyperref[def:mono-epi-morphism]{monomorphism}.

          If \(g\colon Z \to X\) is a morphism such that \(f \circ g = 0\) then by \autoref{lemma:mono-epi-additive} \(g = 0\).
          In particular \(g\) factors uniquely through \(0 \to X\), i.e: \(0 \to X\) is the \hyperref[def:categorical-kernel]{kernel} of \(f\).
    \item[\(\impliedby\)] Suppose now that \(0 \to X\) is the \hyperref[def:categorical-kernel]{kernel} of \(f\).

          Consider a morphism \(g\colon Z \to X\) such that \(f \circ g = 0\).
          By the \hyperref[def:categorical-kernel]{kernel} properties we have the following diagram
          \begin{center}
            \begin{tikzcd}
              Z \arrow[rr, "g"'] \arrow[rrr, bend left, "0"] \arrow[dr, dotted, "\exists!"'] & & X \arrow[r, "f"'] & Y \\
              & 0 \arrow[ur]
            \end{tikzcd}
          \end{center}
          Therefore by uniqueness \(g = 0\), making \(f\) a \hyperref[def:mono-epi-morphism]{monomorphism} by \autoref{lemma:mono-epi-additive}. \qedhere
  \end{itemize}
\end{proof}

As in the case of \hyperref[def:additive-category]{additive categories}, \hyperref[def:functor-category]{functor categories} play nicely with \hyperref[def:abelian-category]{abelian categories}.

\begin{proposition}\label{prop:functor-category-is-abelian}
  Let \(\category{I}\) be a \hyperref[def:small-category]{small category} and \(\category{C}\) an \hyperref[def:abelian-category]{abelian category}.
  Then the \hyperref[def:functor-category]{functor category} \(\Fct{\category{I}}{\category{C}}\) is \hyperref[def:abelian-category]{abelian}.
\end{proposition}
\begin{proof}
  By \autoref{lemma:functor-category-is-additive} we have that \(\Fct{\category{I}}{\category{C}}\) is \hyperref[def:additive-category]{additive}.
  
  Furthermore since \(\category{C}\) has \hyperref[def:categorical-kernel]{kernels} and \hyperref[def:categorical-cokernel]{cokernels}, which are no more than special cases of \hyperref[def:categorical-limit]{limits} and \hyperref[def:categorical-colimit]{colimits}, by \autoref{prop:limits-in-functor-category} \(\Fct{\category{I}}{\category{C}}\) has \hyperref[def:categorical-kernel]{kernels} and \hyperref[def:categorical-cokernel]{cokernels} and they are computed elementwise.

  Finally since \(\operatorname{Coim} f(i) \isoarrow \operatorname{Im} f(i)\) for all \(i \in \category{I}\) we have that \(\operatorname{Coim} f \isoarrow \operatorname{Im} f\).
\end{proof}

Furthermore since \hyperref[def:categorical-kernel]{kernel}, \hyperref[def:categorical-cokernel]{cokernels} and \hyperref[def:direct-sum]{direct sums} are computed term-wise we directly obtain that \(\Complex[\ast]{\category{C}}\) is \hyperref[def:abelian-category]{abelian}.
\begin{corollary}\label{cor:complex-category-is-abelian}
  If \(\category{C}\) is \hyperref[def:abelian-category]{abelian} then the \hyperref[def:complex]{categories of complexes} \(\Complex[\ast]{\category{C}}\) (\(\ast = +, -, \operatorname{b}\)) are \hyperref[def:abelian-category]{abelian}.
\end{corollary}

\section{Exactness}

In what remains of the chapter we will assume that \(\category{C}\) is an \hyperref[def:abelian-category]{abelian category}.


\begin{definition}\label{def:exact-sequence}~
\begin{enumerate}
  \item A \hyperref[def:complex]{complex} \(X \xrightarrow{f} Y \xrightarrow{g} Z\) is \textbf{exact} if \(\operatorname{Im} f \isoarrow \ker g\).
        Note that this is equivalent to \(\coker f \isoarrow \operatorname{Im} g\).
  \item More generally, a sequence of morphisms \(X^{p} \xrightarrow{\d^{p}} \cdots \to X^{n}\) with \(\d^{i+1} \circ \d^{i} = 0\) for all \(i \in \{p, \ldots, n-1\}\) is \textbf{exact} if \(\operatorname{Im} \d^{i} \isoarrow \ker \d^{i+1}\) for all \(i \in \{p, \ldots, n-1\}\).
  \item A \textbf{short exact sequence} is an exact sequence of the form
        \begin{equation*}
          0 \to X \to Y \to Z \to 0
        \end{equation*}
\end{enumerate}
\end{definition}

\begin{remark}
  A short sequence \(0 \to X \xrightarrow{f} Y \xrightarrow{g} Z \to 0\) is \hyperref[def:exact-sequence]{short exact} if and only if \(f\) is a \hyperref[def:mono-epi-morphism]{monomorphism}, \(g\) is an \hyperref[def:mono-epi-morphism]{epimorphism} and \(\operatorname{Im} f = \ker g\).

  Indeed, it is \hyperref[def:exact-sequence]{exact} at \(X\) iff \(\ker f = \operatorname{Im} 0 = 0\) and is \hyperref[def:exact-sequence]{exact} at \(Z\) iff \(\coker g = \operatorname{Im} 0 = 0\).
\end{remark}

\begin{example}
  Any morphism \(f: X \to Y\) may be decomposed in \hyperref[def:exact-sequence]{short exact sequences}
  \begin{center}
    \begin{tikzcd}
      0 \arrow[r] & \ker f \arrow[r, tail] & X \arrow[r, two heads] & \operatorname{Coim} f \arrow[r] & 0 \\
      0 \arrow[r] & \operatorname{Im} f \arrow[r, tail] & Y \arrow[r, two heads] & \coker f \arrow[r] & 0
    \end{tikzcd}
  \end{center}
  with \(\operatorname{Coim} f \cong \operatorname{Im} f\).

  Indeed, by definition we have \(\operatorname{Coim} f = \coker \left( \ker f \to X \right)\) and \(\operatorname{Im} = \ker (Y \to \coker f)\).
\end{example}

A stronger condition than \hyperref[def:exact-sequence]{exactness} is split exactness.
Roughly speaking, a \hyperref[def:exact-sequence]{short exact sequence} is split exact if the only information is at its ends.
\begin{proposition}\label{prop:split-exact-sequence}
  Let
  \begin{center}
    \begin{tikzcd}
      0 \arrow[r] & X \arrow[r, "f"] & Y \arrow[r, "g"] & Z \arrow[r] & 0
    \end{tikzcd}
  \end{center}
  be a \hyperref[def:exact-sequence]{short exact sequence} in an \hyperref[def:abelian-category]{abelian category} \(\category{C}\).
  We say that the sequence \textbf{splits} if any of the following equivalent conditions are satisfied
  \begin{enumerate}
    \item There exists \(h\colon Z \to Y\) such that \(g \circ h = \id_{Z}\).
    \item There exists \(k\colon Y \to X\) such that \(k \circ f = \id_{X}\).
    \item There exists \(\varphi = (k, h)\) and \(\psi = {f \choose g}\) such that \(Y \xrightarrow{\varphi} X \oplus Z\) and \(X \oplus Z \xrightarrow{\psi} Y\) are inverses to each other.
    \item The sequence is \hyperref[def:homotopy]{homotopic to \(0\)}.
    \item The sequence is \hyperref[def:categorical-isomorphism]{isomorphic} to \(0 \to X \to X \oplus Z \to Z \to 0\).
  \end{enumerate}
\end{proposition}
\begin{proof}~
  \begin{enumerate}
    \item[\(1 \implies 3\)] Since \(g = g \circ h \circ g\) we have that \(g \circ (\id_{Y} - h \circ g) = 0\), i.e: \(\id_{Y} - h \circ g\) factors through \(\ker g\) and so it factors through \(X\).
          Therefore there exists \(h\colon Y \to Z\) such that \(\id_{Y} - h \circ g = f \circ k\).
    \item[\(3 \implies 1\)] Since \(g \circ f = 0\) we have \(g = g \circ h \circ g\), i.e: \((g \circ h - \id_{Z}) \circ g\).
          Therefore since \(g\) is an \hyperref[def:mono-epi-morphism]{epimorphism} we have \(g \circ h - \id_{Z} = 0 \iff g \circ h = \id_{Z}\).
    \item[\(2 \iff 3\)] Analogous to \(1 \iff 3\).
    \item[\(3 \iff 4\)] By definition the sequence is \hyperref[def:homotopy]{homopic to zero} if and only if there exists
          \begin{center}
            \begin{tikzcd}
              0 \arrow[r] & X \arrow[d, "\id_{X}"'] \arrow[r, "f"] & Y \arrow[d, "\id_{Y}"'] \arrow[r, "g"] \arrow[dl, dotted, "k"'] & Z \arrow[d, "\id_{Z}"] \arrow[r] \arrow[dl, dotted, "h"'] & 0 \\
              0 \arrow[r] & X \arrow[r, "f"'] & Y \arrow[r, "g"'] & Z \arrow[r] & 0
            \end{tikzcd}
          \end{center}
          such that \(\id_{X} = k \circ f\), \(\id_{Y} = h \circ g\) and \(\id_{Z} = h \circ g + f \circ k\).
    \item[\(3 \iff 5\)] Trivial. \qedhere
  \end{enumerate}
\end{proof}

When studying \hyperref[def:exact-sequence]{exactness} between different \hyperref[def:category]{categories} we don't want to lose information on whether a sequence is \hyperref[def:exact-sequence]{exact} or not.
Therefore the correct notion of \hyperref[def:functor]{functors} in this case is that of exact functors.
\begin{definition}\label{def:exact-functor}
  Let \(\category{C}, \category{D}\) be two \hyperref[def:abelian-category]{abelian categories} and
  \begin{center}
    \begin{tikzcd}
      0 \arrow[r] & X \arrow[r] & Y \arrow[r] & Z \arrow[r] & 0
    \end{tikzcd}
  \end{center}
  be a \hyperref[def:exact-sequence]{short exact sequence} in an \(\category{C}\).
  
  An \hyperref[def:additive-functor]{additive functor} \(F\colon \category{C} \to \category{D}\) is \textbf{left-exact} if
   \begin{center}
    \begin{tikzcd}
      0 \arrow[r] & F(X) \arrow[r] & F(Y) \arrow[r] & F(Z)
    \end{tikzcd}
  \end{center}
  is \hyperref[def:exact-sequence]{exact} in \(\category{D}\).

  Dually, it is \textbf{right-exact} if
   \begin{center}
    \begin{tikzcd}
      F(X) \arrow[r] & F(Y) \arrow[r] & F(Z) \arrow[r] & 0
    \end{tikzcd}
  \end{center}
  is \hyperref[def:exact-sequence]{exact} in \(\category{D}\).

  If \(F\) is both left and right exact we say that it is \textbf{exact}.
\end{definition}

We state the following characterization of \hyperref[def:exact-functor]{exact functors}, which we will not prove.
\begin{lemma}[{\cite[Lemma 5.2.1]{kashiwaraIntroductionCategoriesHomological}}]\label{lemma:characterization-exact-functor}
  Let \(F\colon \category{C} \to \category{D}\) be an \hyperref[def:additive-functor]{additive functor}.
  \begin{enumerate}
    \item The following are equivalent
          \begin{enumerate}
            \item \(F\) is \hyperref[def:exact-functor]{left exact} (resp. \hyperref[def:exact-functor]{right exact}).
            \item \(F\) commutes with \hyperref[def:categorical-kernel]{kernels}: for any morphism \(f\colon X \to Y\), \(F(\ker f) \isoarrow \ker Ff\) (resp. with \hyperref[def:categorical-cokernel]{cokernels}).
            \item \(F\) commutes with \hyperref[def:categorical-limit]{finite limits} (resp. \hyperref[def:categorical-colimit]{finite colimits}).
            \item For any \hyperref[def:exact-sequence]{short exact sequence} \(0 \to X \to Y \to Z\) (resp. \(X \to Y \to Z \to 0\)) in \(\category{C}\) the sequence \(0 \to F(X) \to F(Y) \to F(Z)\) is \hyperref[def:exact-sequence]{exact} in \(\category{D}\) (resp. \(F(X) \to F(Y) \to F(Z) \to 0\) is \hyperref[def:exact-sequence]{exact}).
          \end{enumerate}
    \item  The following are equivalent
          \begin{enumerate}
            \item \(F\) is \hyperref[def:exact-functor]{exact}.
            \item For any \hyperref[def:exact-sequence]{exact sequence} \(X \to Y \to Z\) in \(\category{C}\) the sequence \(F(X) \to F(Y) \to F(Z)\) is \hyperref[def:exact-sequence]{exact} in \(\category{D}\).
            \item \(F\) commutes with \hyperref[def:categorical-colimit]{finite limits} and \hyperref[def:categorical-colimit]{finite colimits}.
          \end{enumerate}
  \end{enumerate}
\end{lemma}

As a direct application of \autoref{lemma:characterization-exact-functor} we have the following examples.
\begin{example}\label{ex:hom-left-exact}
  The \hyperref[ex:hom-functor]{Hom bifunctor} \(\Hom_{\category{C}}\colon \opcategory{C} \times \category{C} \to \Ab\) is \hyperref[def:exact-functor]{left exact} with respect to each of its arguments.
\end{example}

\begin{example}\label{ex:adjoint-functors-and-exactness}
  Recall \autoref{thm:adjoints-preserve-limits}.
  
  \hyperref[def:adjoint-functors]{Right adjoints} are \hyperref[def:exact-functor]{left exact}.
  Dually, \hyperref[def:adjoint-functors]{left adjoints} are \hyperref[def:exact-functor]{right exact}.
\end{example}

\begin{example}\label{ex:limit-colimit-exact}
  Let \(\category{I}\) be a \hyperref[def:small-category]{small category}.
  If \(\category{C}\) admits \hyperref[def:categorical-limit]{\(I\)-indexed limits} then the \hyperref[def:functor]{functor} \(\lim\colon \Fct{\opcategory{I}}{\category{C}} \to \category{C}\) is \hyperref[def:exact-functor]{left exact}.
  Dually if \(\category{C}\) admits \hyperref[def:categorical-colimit]{\(I\)-indexed colimits} then the \hyperref[def:functor]{functor} \(\colim\colon \Fct{\category{I}}{\category{C}} \to \category{C}\) is \hyperref[def:exact-functor]{right exact}.
\end{example}

% \begin{example}\label{ex:product-direct-sum-exact}
%   Let \(R\) be a ring and \(\mathcal{I}\) an index set.
%   The two \hyperref[def:functor]{functors}
%   \begin{equation*}
%     \prod_{i \in \mathcal{I}}\colon \Fct{\mathcal{I}}{\Mod(R)} \to \Mod(R) \qquad \bigoplus_{i \in \mathcal{I}} :\Fct{\mathcal{I}}{\Mod(R)} \to \Mod(R)
%   \end{equation*}
%   are \hyperref[def:exact-functor]{exact}.
% \end{example}

% \begin{example}\label{ex:injlim-exact}
%   Let \(R\) be a ring and \(\mathcal{I}\) a \hyperref[def:directed-set]{small directed category}.
%   The \hyperref[def:functor]{functor} \(\varinjlim_{i \in \mathcal{I}}\colon \Fct{\mathcal{I}}{\Mod(R)} \to \Mod(R)\) is \hyperref[def:exact-functor]{exact}.
% \end{example}

\subsection{Diagram lemmas}

Note that \hyperref[def:abelian-category]{abelian categories} may not necessarily consists of sets with elements (for example the \hyperref[def:complex]{category of complexes}).
Therefore proving facts about these \hyperref[def:category]{categories} requires, in principle, to work over the full abstraction of category theory.
However using the \hyperref[thm:freyd-mitchell]{Freyd-Mitchell embedding theorem} one can \textit{cheat} a little.

\begin{theorem}[{Freyd-Mitchell embedding - \cite[Section 1.6]{weibelIntroductionHomologicalAlgebra1994}}]\label{thm:freyd-mitchell}
  Let \(\category{C}\) be a \hyperref[def:small-category]{small} \hyperref[def:abelian-category]{abelian category}.
  Then there exists a ring \(R\) (with unit, not necessarilly commutative) and a \hyperref[def:functor-types]{fully faithful} \hyperref[def:exact-functor]{exact} \hyperref[def:functor]{functor} \(F\colon \category{C} \to \Mod(R)\) which embeds \(\category{C}\) as a \hyperref[def:subcategory]{full subcategory} of \(\Mod(R)\).
\end{theorem}

In short, to prove something about a \hyperref[def:diagram]{diagram} in an \hyperref[def:abelian-category]{abelian category} we may assume that we are working in the \hyperref[def:category]{category} of modules over a ring \(R\) and then diagram-chase elements.
Furthermore, any fact about \hyperref[def:categorical-kernel]{kernels}, \hyperref[def:categorical-cokernel]{cokernels} and so on that holds in \(\Mod(R)\) also holds in any \hyperref[def:abelian-category]{abelian category}.

% Another strategy to prove statements in \hyperref[def:abelian-category]{abelian categories} is the \textit{philosophy of generalized elements} \cite[Theorem VIII.4.3]{maclaneCategoriesWorkingMathematician1978}.

Using \hyperref[thm:freyd-mitchell]{this theorem} we can easilly prove two diagram lemmas that will make our lives easier later on.

\begin{lemma}[Five lemma]\label{lemma:five-lemma}
  Consider a commutative diagram
  \begin{center}
    \begin{tikzcd}
      X^{0} \arrow[r, "\alpha^{0}"] \arrow[d, "f^{0}"'] & X^{1} \arrow[r, "\alpha^{1}"] \arrow[d, "f^{1}"] & X^{2} \arrow[r, "\alpha^{2}"] \arrow[d, "f^{2}"] & X^{3} \arrow[d, "f^{3}"] \\
      Y^{0} \arrow[r, "\beta^{0}"'] & Y^{1} \arrow[r, "\beta^{1}"'] & Y^{2} \arrow[r, "\beta^{2}"'] & Y^{3}
    \end{tikzcd}
  \end{center}
  where the rows are \hyperref[def:exact-sequence]{exact}.

  If \(f^{0}\) is an \hyperref[def:mono-epi-morphism]{epimorphism} and \(f^{1}, f^{3}\) are \hyperref[def:mono-epi-morphism]{monomorphisms} then \(f^{2}\) is a \hyperref[def:mono-epi-morphism]{monomorphism}.
  Dually if \(f^{3}\) is an \hyperref[def:mono-epi-morphism]{monomorphism} and \(f^{0}, f^{2}\) are \hyperref[def:mono-epi-morphism]{epimorphisms} then \(f^{1}\) is a \hyperref[def:mono-epi-morphism]{epimorphism}.
\end{lemma}
\begin{proof}
  Let \(x_{2} \in X^{2}\) such that \(f^{2}(x_{2}) = 0\).
  We want to see that \(x_{2} = 0\).
  
  We have that \((f^{3} \circ \alpha^{2})(x_{2}) = 0\) which, since \(f^{3}\) is a \hyperref[def:mono-epi-morphism]{monomorphism}, implies that \(\alpha^{2}(x_{2}) = 0\).
  Since the first row is \hyperref[def:exact-sequence]{exact} there exists \(x_{1} \in X^{1}\) such that \(\alpha(x_{1}) = x_{2}\).

  Let \(y_{1} \coloneq f^{1}(x_{1})\).
  Since \((\beta^{1} \circ f^{1})(x_{1}) = 0\) and the second row is \hyperref[def:exact-sequence]{exact} there exists \(y_{0} \in X^{0}\) such that \(\beta^{0}(y_{0}) = f^{1}(x_{1})\).

  Since \(f^{0}\) is an \hyperref[def:mono-epi-morphism]{epimorphism} there exists \(x_{0} \in X^{0}\) such that \(y_{0} = f^{0}(x_{0})\).
  
  Since \((f^{1} \circ \alpha^{0})(x_{0}) = f^{1}(x_{1})\) and \(f^{1}\) is a \hyperref[def:mono-epi-morphism]{monomorphism} we have \(\alpha^{0}(x_{0}) = x_{1}\).
  Therefore \(x_{2} = \alpha^{1}(x_{1}) = 0\).
\end{proof}


\begin{lemma}[Snake lemma]\label{lemma:snake-lemma}
  Consider a commutative diagram
  \begin{center}
    \begin{tikzcd}
      & X' \arrow[r, "f"] \arrow[d, "\alpha"'] & X \arrow[r, two heads, "g"] \arrow[d, "\beta"] & X'' \arrow[r] \arrow[d, "\gamma"] & 0\\
      0 \arrow[r] & Y' \arrow[r, tail, "f'"'] & Y \arrow[r, "g'"'] & Y''
    \end{tikzcd}
  \end{center}
  where the rows are \hyperref[def:exact-sequence]{exact}.

  Then there exists a morphism \(\delta\colon \ker \gamma \to \coker \alpha\) giving rise to an \hyperref[def:exact-sequence]{exact sequence}
  \begin{center}
    \begin{tikzcd}
        & 0 \arrow[d] & 0 \arrow[d] & 0 \arrow[d]  &   \\
        & \ker \alpha \arrow[d, tail] \arrow[r, dashed] & \ker \beta \arrow[d, tail] \arrow[r, dashed] \arrow[ddd, phantom, ""{coordinate, name=Z}] & \ker \gamma \arrow[d, tail] \arrow[dddll, dashed, rounded corners, to path={ -- ([xshift=2ex]\tikztostart.east)
          |- (Z) [near end]\tikztonodes
          -| ([xshift=-2ex]\tikztotarget.west)
          -- (\tikztotarget)}] &   \\
        & X' \arrow[d, "\alpha" near start] \arrow[r, "f"] & X \arrow[d, "\beta" near start] \arrow[r, two heads, "g"] & X'' \arrow[d, "\gamma" near start] \arrow[r] & 0 \\
      0 \arrow[r] & Y' \arrow[d, two heads] \arrow[r, tail, "f'"'] & Y \arrow[d, two heads] \arrow[r, "g'"'] & Y'' \arrow[d, two heads] &   \\
        & \coker \alpha \arrow[d] \arrow[r, dashed] & \coker \beta \arrow[d] \arrow[r, dashed] & \coker \gamma \arrow[d] &   \\
        & 0  & 0 & 0
    \end{tikzcd}
  \end{center}
\end{lemma}
\begin{proof}
  We will first check that \(\ker \alpha \to \ker \beta \to \ker \gamma\) is \hyperref[def:exact-sequence]{exact}.
  Let \(x \in \ker \beta\) with \(g(x) = 0\).
  Since the first row is \hyperref[def:exact-sequence]{exact} there exists \(x' \in X'\) with \(f(x') = x\).
  Then \((f' \circ \alpha)(x') = (\beta \circ f)(x') = 0\).
  Since \(f'\) is a \hyperref[def:mono-epi-morphism]{monomorphism} we get \(\alpha(x') = 0\) and thus \(x' \in \ker \alpha\).

  A dual argument shows that \(\coker \alpha \to \coker \beta \to \coker \gamma\) is \hyperref[def:exact-sequence]{exact}.

  It only remains to construct \(\delta\).
  Let \(x'' \in \ker \gamma\) and choose \(x \in X\) with \(g(x) = x''\).
  Set \(y \coloneq \beta(y)\).
  Since \(g'(y) = 0\) there exists \(y' \in Y'\) such that \(f'(y') = y\).
  Therefore we can define \(\delta(x'')\) to be the image of \(y'\) in \(\coker \alpha\).
\end{proof}


\section{Complexes and cohomology}

Recall that \hyperref[cor:complex-category-is-abelian]{the categories \(\Complex[\ast]{\category{C}}\) are abelian for \(\ast = +, -, \operatorname{b}, \operatorname{ub}\)}.

Since for a \hyperref[def:complex]{complex} \(X^{\bullet}\) we always have \(\operatorname{Im} \d_{X}^{n} \subseteq \ker \d_{X}^{n-1}\) for all \(n \in \mathbb{N}\), there is a failure of \hyperref[def:exact-sequence]{exactness} if and only if \(\operatorname{Im} \d_{X}^{n} \neq \ker \d_{X}^{n-1}\).
We introduce cohomology as a way of measuring how far away is a \hyperref[def:complex]{complex} from being \hyperref[def:exact-sequence]{exact}.

\begin{definition}[Cohomology]\label{def:cohomology}
  Let
  \begin{center}
    \begin{tikzcd}
      X^{\bullet} \coloneq \cdots \arrow[r] & X^{n-1} \arrow[r, "\d_{X}^{n-1}"] & X^{n} \arrow[r, "\d_{X}^{n}"] & X^{n+1} \arrow[r] & \cdots
    \end{tikzcd}
  \end{center}
  be a \hyperref[def:complex]{complex} in \(\category{C}\) (i.e: \(X^{\bullet} \in \Complex{\category{C}}\)).
  We define
  \begin{align*}
    Z^{n}(X^{\bullet}) &\coloneq \ker \d_{X}^{n}  \\
    B^{n}(X^{\bullet}) &= \operatorname{Im} \d_{X}^{n-1} \\
    \cohomology{n}(X^{\bullet}) &\coloneq \quot{Z^{n}(X^{\bullet})}{B^{n}(X^{\bullet})} = \coker(B^{n}(X^{\bullet}) \to Z^{n}(X^{\bullet}))
  \end{align*}
  We call \(\cohomology{n}(X^{\bullet})\) the \(n\)-th \textbf{cohomology} of \(X^{\bullet}\).

  Note that this defines a \hyperref[def:functor]{functor} \(\cohomology{n}\colon \Complex{\category{C}} \to \category{C}\).
\end{definition}

\begin{remark}
  A \hyperref[def:complex]{complex} \(X^{\bullet}\) is \hyperref[def:exact-sequence]{exact} at \(X^{n}\) if and only if \(\cohomology{i}(X^{\bullet}) = 0\).
\end{remark}

\begin{remark}\label{rem:short-exact-sequence-cohomology}
  We have \hyperref[def:exact-sequence]{short exact sequences}
  \begin{center}
    \begin{tikzcd}
      X^{n-1} \arrow[r, "\d_{X}^{n-1}"] & \ker \d_{X}^{n} \arrow[r] & \cohomology{n}(X^{\bullet}) \arrow[r] & 0 \\
      0 \arrow[r] & \cohomology{n}(X^{\bullet}) \arrow[r] & \coker \d_{X}^{n-1} \arrow[r, "\d_{X}^{n}"] & X^{n+1}
    \end{tikzcd}
  \end{center}

  We can glue this into an \hyperref[def:exact-sequence]{exact sequence}
  \begin{center}
    \begin{tikzcd}
      0 \arrow[r] & \cohomology{n}(X^{\bullet}) \arrow[r] & \coker \d_{X}^{n-1} \arrow[r, "\d_{X}^{n}"] & \ker \d_{X}^{n+1} \arrow[r] & \cohomology{n+1}(X^{\bullet}) \arrow[r] & 0
    \end{tikzcd}
  \end{center}
\end{remark}

The first thing we are going to see is that \hyperref[def:cohomology]{cohomology} is an \hyperref[def:homotopy]{homotopy} invariant.
\begin{lemma}\label{lemma:extension-cohomology-homotopy}
  Let \(\category{C}\) be an \hyperref[def:abelian-category]{abelian category} and \(f^{\bullet}\colon X^{\bullet} \to Y^{\bullet}\) a morphism in \(\Complex{\category{C}}\) \hyperref[def:homotopy]{homotopic to zero}.
  Then \(\cohomology{n}f^{\bullet}\colon \cohomology{n}(X^{\bullet}) \to \cohomology{n}(Y^{\bullet})\) is the \hyperref[def:zero-object]{zero morphism} for all \(n \in \mathbb{Z}\).

  In particular if \(f^{\bullet}, g^{\bullet}: X^{\bullet} \to Y^{\bullet}\) are \hyperref[def:homotopy]{homotopic} \(\cohomology{n}f^{\bullet} = \cohomology{n}g^{\bullet}\) for all \(n \in \mathbb{Z}\).
\end{lemma}
\begin{proof}
  Let
  \begin{equation*}
    f^{n} = s^{n+1} \circ \d_{X}^{n} + \d_{Y}^{n-1} \circ s^{n}
  \end{equation*}
  Then \(\d_{X}^{n} = 0\) in \(\ker \d_{X}^{n}\) and the image of \(\d_{Y}^{n-1}\) is zero on \(\quot{\ker \d_{Y}^{n}}{\operatorname{Im} \d_{Y}^{n-1}}\).
  
  Therefore \(\cohomology{n}f^{\bullet}: \quot{Z^{n}(X^{\bullet})}{B^{n}(X^{\bullet})} \to \quot{Z^{n}(Y^{\bullet})}{B^{n}(Y^{\bullet})}\) is the zero morphism.
\end{proof}

\begin{corollary}\label{cor:homotopy-preserves-cohomology}
  Let \(X^{\bullet} \sim Y^{\bullet}\) be two \hyperref[def:homotopy]{homotopic complexes}.
  Then \(\cohomology{n}(X^{\bullet}) \cong \cohomology{n}(Y^{\bullet})\).
\end{corollary}
\begin{proof}
  Since \(X^{\bullet} \sim Y^{\bullet}\) there exist maps \(f^{\bullet}: X^{\bullet} \to Y^{\bullet}\) and \(g^{\bullet}: Y^{\bullet} \to X^{\bullet}\) such that \(g^{\bullet} \circ f^{\bullet} \sim \id_{X^{\bullet}}\) and \(f^{\bullet} \circ g^{\bullet} \sim \id_{Y^{\bullet}}\).

  Therefore by \autoref{lemma:extension-cohomology-homotopy} for all \(n \in \mathbb{Z}\) we have that
  \begin{align*}
    \cohomology{n}(g^{\bullet} \circ f^{\bullet}) &\cong \cohomology{n} \id_{X^{\bullet}} = \id_{\cohomology{n}(X^{\bullet})} \\
    \cohomology{n}(f^{\bullet} \circ g^{\bullet}) &\cong \cohomology{n} \id_{Y^{\bullet}} = \id_{\cohomology{n}(Y^{\bullet})} 
  \end{align*}
  Furthermore, since \(\cohomology{n}\) is a \hyperref[def:functor]{functor}
  \begin{align*}
    \cohomology{n}(g^{\bullet} \circ f^{\bullet}) &= \cohomology{n}g^{\bullet} \circ \cohomology{n}f^{\bullet} \\
    \cohomology{n}(f^{\bullet} \circ g^{\bullet}) &= \cohomology{n}f^{\bullet} \circ \cohomology{n}g^{\bullet} 
  \end{align*}
  which means that \(\cohomology{n}f^{\bullet}\) and \(\cohomology{n}g^{\bullet}\) are inverses of each other.
\end{proof}

Since \hyperref[def:cohomology]{cohomology is a functor} we can also study how it acts on morphisms.
For some of these morphisms, the resulting \hyperref[def:cohomology]{cohomology objects} will be the same and thus they are ``equivalent'' in some sense.
This leads to the notion of a quasi-isomorphism.
\begin{definition}[Quasi-isomorphism]\label{def:quasi-isomorphism}
  Let \(\category{C}\) be an \hyperref[def:abelian-category]{abelian category}.
  A morphism \(f^{\bullet}\colon X^{\bullet} \to Y^{\bullet}\) is a \textbf{quasi-isomorphism} (abbreviated q.is) if \(\cohomology{n}f^{\bullet}\) is an \hyperref[def:categorical-isomorphism]{isomorphism} for all \(n \in \mathbb{Z}\).
  In this case we say that \(X^{\bullet}\) and \(Y^{\bullet}\) are \textbf{quasi-isomorphic}.

  In particular \(X^{\bullet}\) is quasi-isomorphic to zero if and only if \(X^{\bullet}\) is \hyperref[def:exact-sequence]{exact}.
\end{definition}

Recall that we can identify an object \(X \in \category{C}\) with a complex \(\cdots \to  0 \to X \to 0 \to \cdots\).
Using \hyperref[def:quasi-isomorphism]{quasi-isomorphisms} we can associated this bounded complex an associated unbounded complex that has the same \hyperref[def:cohomology]{cohomology}.
\begin{definition}[Resolution]\label{def:resolution}
  Let \(\category{C}\) be an \hyperref[def:abelian-category]{abelian category} and \(X \in \category{C}\).
  A \textbf{right resolution} of \(X\) is a \hyperref[def:quasi-isomorphism]{quasi-isomorphism} \(X \to Y^{\bullet}\) with \(Y^{\bullet} \in \Complex[+]{\category{C}}\).

  Dually, a \textbf{left resolution} is a \hyperref[def:quasi-isomorphism]{quasi-isomorphism} \(Z^{\bullet} \to X\) with \(Z^{\bullet} \in \Complex[-]{\category{C}}\).
\end{definition}

\begin{remark}
  Let \(X \to Y^{\bullet}\) be a \hyperref[def:resolution]{right resolution} of \(X\).

  This means that we have a \hyperref[def:complex]{morphism of complexes}
  \begin{center}
    \begin{tikzcd}
      \cdots \arrow[r] & 0 \arrow[r] \arrow[d] & X \arrow[r] \arrow[d, "\varepsilon"'] & 0 \arrow[r] \arrow[d] & \cdots \\
      \cdots \arrow[r] & 0 \arrow[r] & Y^{0} \arrow[r] & Y^{1} \arrow[r] & \cdots
    \end{tikzcd}
  \end{center}
  Note that this is exactly the same data as an \hyperref[def:exact-sequence]{exact complex}
  \begin{center}
    \begin{tikzcd}
      0 \arrow[r] & X \arrow[r, "\varepsilon"] & Y^{0} \arrow[r] & Y^{1} \arrow[r] & \cdots
    \end{tikzcd}
  \end{center}
\end{remark}

We end the chapter with a key result: taking \hyperref[def:cohomology]{cohomology} transforms a \hyperref[def:exact-sequence]{short exact sequence} of \hyperref[def:complex]{complexes} into a \hyperref[def:exact-sequence]{long exact sequence} in a functorial way.
\begin{theorem}[Long exact sequence in cohomology]\label{thm:LES-in-cohomology}
  Let
  \begin{center}
    \begin{tikzcd}
      0 \arrow[r] & X^{\bullet} \arrow[r] & Y^{\bullet} \arrow[r] & Z^{\bullet} \arrow[r] & 0
    \end{tikzcd}
  \end{center}
  be a \hyperref[def:exact-sequence]{short exact sequence} in an \hyperref[def:abelian-category]{abelian category} \(\Complex{\category{C}}\).
  \begin{enumerate}
    \item For all \(k \in \mathbb{Z}\) the sequence \(\cohomology{k}(X^{\bullet}) \to \cohomology{k}(Y^{\bullet}) \to \cohomology{k}(Z^{\bullet})\) is \hyperref[def:exact-sequence]{exact}.
    \item For all \(k \in \mathbb{Z}\) there exists \(\delta^{k}\colon \cohomology{k}(Z^{\bullet}) \to \cohomology{k+1}(X^{\bullet})\) making the long sequence
          \begin{center}
            \begin{tikzcd}
              \cdots \arrow[r] & \cohomology{k-1}(Y^{\bullet}) \arrow[r] \arrow[d, phantom, ""{coordinate, name=X}] & \cohomology{k-1}(Z^{\bullet}) \arrow[dll, rounded corners, "\delta^{k-1}" description, to path={
                -- ([xshift=2ex]\tikztostart.east)
                |- (X) [near end]\tikztonodes
                -| ([xshift=-2ex]\tikztotarget.west)
                -- (\tikztotarget)}] \\
              \cohomology{k}(X^{\bullet}) \arrow[r] & \cohomology{k}(Y^{\bullet}) \arrow[r] \arrow[d, phantom, ""{coordinate, name=Y}] & \cohomology{k}(Z) \arrow[dll, rounded corners, "\delta^{k}" description, to path={
                -- ([xshift=2ex]\tikztostart.east)
                |- (Y) [near end]\tikztonodes
                -| ([xshift=-2ex]\tikztotarget.west)
                -- (\tikztotarget)}] \\
              \cohomology{k+1}(X^{\bullet}) \arrow[r] & \cohomology{k+1}(Y^{\bullet}) \arrow[r] & \cdots
            \end{tikzcd}
          \end{center}
          \hyperref[def:exact-sequence]{exact}.
          Furthermore \(\delta^{k}\) can be constructed to be functorial with respect to \hyperref[def:exact-sequence]{short exact sequences} of \(\Complex{\category{C}}\).
  \end{enumerate}
\end{theorem}
\begin{proof}
  Consider the diagram
  \begin{center}
    \begin{tikzcd}
      & 0 \arrow[d] & 0 \arrow[d] & 0 \arrow[d] \\
      & \cohomology{k}(X^{\bullet}) \arrow[d] & \cohomology{k}(Y^{\bullet}) \arrow[d] & \cohomology{k}(Z^{\bullet}) \arrow[d] \\
      & \coker \d_{X}^{k-1} \arrow[r] \arrow[d] & \coker \d_{Y}^{k-1} \arrow[r] \arrow[d] & \coker \d_{Z}^{k-1} \arrow[r] \arrow[d] & 0 \\
      0 \arrow[r] & \ker \d_{X}^{k+1} \arrow[r] \arrow[d] & \ker \d_{Y}^{k+1} \arrow[r] \arrow[d] & \ker \d_{Z}^{k+1} \arrow[d] \\
      & \cohomology{k+1}(X^{\bullet}) \arrow[d] & \cohomology{k+1}(Y^{\bullet}) \arrow[d] & \cohomology{k+1}(Z^{\bullet}) \arrow[d] \\
      & 0 & 0 & 0 
    \end{tikzcd}
  \end{center}
  where the columns are \hyperref[def:exact-sequence]{exact} by \autoref{rem:short-exact-sequence-cohomology} and the rows are \hyperref[def:exact-sequence]{exact} by hypothesis.
  Therefore the result follows from the \hyperref[lemma:snake-lemma]{Snake lemma}.
\end{proof}

% \begin{corollary}
%   For any morphism \(f^{\bullet}: X^{\bullet} \to Y^{\bullet}\) in \(\Complex{\category{C}}\) there exists a \hyperref[def:exact-sequence]{long exact sequence}
%   \begin{center}
%     \begin{tikzcd}
%       \cdots \arrow[r] & \cohomology{k-1}(Y^{\bullet}) \arrow[r] \arrow[d, phantom, ""{coordinate, name=X}] & \cohomology{k-1}(\operatorname{Cone}(f^{\bullet}) \arrow[dll, rounded corners, "\delta^{k-1}" description, to path={
%         -- ([xshift=2ex]\tikztostart.east)
%         |- (X) [near end]\tikztonodes
%         -| ([xshift=-2ex]\tikztotarget.west)
%         -- (\tikztotarget)}] \\
%       \cohomology{k}(X^{\bullet}) \arrow[r] & \cohomology{k}(Y^{\bullet}) \arrow[r] \arrow[d, phantom, ""{coordinate, name=Y}] & \cohomology{k}(\operatorname{Cone}(f^{\bullet})) \arrow[dll, rounded corners, "\delta^{k}" description, to path={
%         -- ([xshift=2ex]\tikztostart.east)
%         |- (Y) [near end]\tikztonodes
%         -| ([xshift=-2ex]\tikztotarget.west)
%         -- (\tikztotarget)}] \\
%       \cohomology{k+1}(X^{\bullet}) \arrow[r] & \cohomology{k+1}(Y^{\bullet}) \arrow[r] & \cdots
%     \end{tikzcd}
%   \end{center}
% \end{corollary}
% \begin{proof}
%   We have a \hyperref[def:complex]{complex}
%   \begin{center}
%     \begin{tikzcd}
%       0 \arrow[r] & Y^{\bullet} \arrow[r] & \operatorname{Cone}(f^{\bullet}) \arrow[r] & X[1]^{\bullet}
%     \end{tikzcd}
%   \end{center}
%   Since in degree \(n\) this gives the \hyperref[prop:split-exact-sequence]{split sequence}
%   \begin{center}
%     \begin{tikzcd}
%       0 \arrow[r] & Y^{n} \arrow[r] & Y^{n} \oplus X^{n+1} \arrow[r] & X^{n+1}
%     \end{tikzcd}
%   \end{center}
%   the \hyperref[def:complex]{complex} is \hyperref[def:exact-sequence]{exact}.
%   Therefore the result follows from \autoref{thm:LES-in-cohomology}.
% \end{proof}

% \subsection{Truncation functors}

% \begin{definition}[Truncation functor]\label{def:truncation-functor}
%   The \textbf{truncation functors}
%   \begin{align*}
%     \tau^{\leq n}, \widetilde{\tau}^{\leq n}:\,& \Complex{\category{C}} \to \Complex[-]{\category{C}} \\
%     \tau^{\geq n}, \widetilde{\tau}^{\geq n}:\,& \Complex{\category{C}} \to \Complex[+]{\category{C}}
%   \end{align*}
%   are defined as, for any \hyperref[def:complex]{complex} \(X^{\bullet} \in \Complex{\category{C}}\)
%   \begin{center}
%     \begin{tikzcd}[column sep=small]
%       \tau^{\leq n}(X^{\bullet}) \coloneq& \cdots \arrow[r] & X^{n-2} \arrow[r] & X^{n-1} \arrow[r] & \ker \d_{X}^{n} \arrow[r] & 0 \arrow[r] & 0 \arrow[r] & \cdots \\
%       \widetilde{\tau}^{\leq n}(X^{\bullet}) \coloneq& \cdots \arrow[r] & X^{n-2} \arrow[r] & X^{n-1} \arrow[r] & X^{n} \arrow[r] & \operatorname{Im} \d_{X}^{n} \arrow[r] & 0 \arrow[r] & \cdots \\
%       \tau^{\geq n}(X^{\bullet}) \coloneq& \cdots \arrow[r] & 0 \arrow[r] & 0 \arrow[r] & \coker \d_{X}^{n-1} \arrow[r] & X^{n+1} \arrow[r] & X^{n+2} \arrow[r] & \cdots \\
%       \widetilde{\tau}^{\geq n}(X^{\bullet}) \coloneq& \cdots \arrow[r] & 0 \arrow[r] & \operatorname{Im} \d_{X}^{n-1} \arrow[r] & X^{n} \arrow[r] & X^{n+1} \arrow[r] & X^{n+2} \arrow[r] & \cdots
%     \end{tikzcd}
%   \end{center}
% \end{definition}

% \begin{remark}\label{rem:truncation-functor-sequences}
%   For any \hyperref[def:complex]{complex} \(X^{\bullet} \in \Complex{\category{C}}\) we have a chain of morphisms
%   \begin{equation*}
%     \tau^{\leq n}(X^{\bullet}) \to \widetilde{\tau}^{\leq n}(X^{\bullet}) \to X^{\bullet} \to \widetilde{\tau}^{\geq n}(X^{\bullet}) \to \tau^{\geq n}(X^{\bullet})
%   \end{equation*}
%   in \(\Complex{\category{C}}\) given by
%   \begin{center}
%     \begin{tikzcd}[column sep=small]
%       \tau^{\leq n}(X^{\bullet}) \arrow[d] =& \cdots \arrow[r] & X^{n-2} \arrow[r] \arrow[d, equal] & X^{n-1} \arrow[r] \arrow[d, equal] & \ker \d_{X}^{n} \arrow[r] \arrow[d] & 0 \arrow[r] \arrow[d] & 0 \arrow[r] \arrow[d, equal] & \cdots \\
%       \widetilde{\tau}^{\leq n}(X^{\bullet}) \arrow[d] =& \cdots \arrow[r] & X^{n-2} \arrow[r] \arrow[d, equal] & X^{n-1} \arrow[r] \arrow[d, equal] & X^{n} \arrow[r] \arrow[d, equal] & \operatorname{Im} \d_{X}^{n} \arrow[r] \arrow[d] & 0 \arrow[r] \arrow[d] & \cdots \\
%       X^{\bullet} \arrow[d] =& \cdots \arrow[r] & X^{n-2} \arrow[r] \arrow[d] & X^{n-1} \arrow[r] \arrow[d] & X^{n} \arrow[r] \arrow[d, equal] & X^{n+1} \arrow[r] \arrow[d, equal] & X^{n+2} \arrow[r] \arrow[d, equal] & \cdots \\
%       \widetilde{\tau}^{\geq n}(X^{\bullet}) \arrow[d] =& \cdots \arrow[r] & 0 \arrow[r] \arrow[d, equal] & \operatorname{Im} \d_{X}^{n-1} \arrow[r] \arrow[d] & X^{n} \arrow[r] \arrow[d] & X^{n+1} \arrow[r] \arrow[d, equal] & X^{n+2} \arrow[r] \arrow[d, equal] & \cdots \\
%       \tau^{\geq n}(X^{\bullet}) =& \cdots \arrow[r] & 0 \arrow[r] & 0 \arrow[r] & \coker \d_{X}^{n-1} \arrow[r] & X^{n+1} \arrow[r] & X^{n+2} \arrow[r] & \cdots \\
%     \end{tikzcd}
%   \end{center}

%   Therefore we get \hyperref[def:exact-sequence]{exact sequences}
%   \begin{center}
%     \begin{tikzcd}
%       0 \arrow[r] & \widetilde{\tau}^{\leq n-1}(X^{\bullet}) \arrow[r] & \tau^{\leq n}(X^{\bullet}) \arrow[r] & \cohomology{n}(X^{\bullet})[-n] \arrow[r] & 0 \\
%       0 \arrow[r] & \cohomology{n}(X^{\bullet})[-n] \arrow[r] & \tau^{\geq n}(X^{\bullet}) \arrow[r] & \widetilde{\tau}^{\geq n+1}(X^{\bullet}) \arrow[r] & 0 \\
%       0 \arrow[r] & \tau^{\leq n}(X^{\bullet}) \arrow[r] & X^{\bullet} \arrow[r] & \widetilde{\tau}^{n+1}(X^{\bullet}) \arrow[r] & 0 \\
%       0 \arrow[r] & \widetilde{\tau}^{\leq n-1}(X^{\bullet}) \arrow[r] & X^{\bullet} \arrow[r] & \tau^{n}(X^{\bullet}) \arrow[r] & 0
%     \end{tikzcd}
%   \end{center}

%   Furthermore we have
%   \begin{align*}
%     \cohomology{j}(\tau^{\leq n}(X^{\bullet})) &\cong \cohomology{j}(\widetilde{\tau}^{\leq n}(X^{\bullet})) = \begin{cases} \cohomology{j}(X^{\bullet}) & j \leq n \\ 0 & j > n \end{cases} \\
%     \cohomology{j}(\tau^{\geq n}(X^{\bullet})) &\cong \cohomology{j}(\widetilde{\tau}^{\geq n}(X^{\bullet})) = \begin{cases} \cohomology{j}(X^{\bullet}) & j \geq n \\ 0 & j < n \end{cases}
%   \end{align*}
% \end{remark}

% \section{Double complexes}

% Recall the \hyperref[rem:finiteness-condition-double-complex]{finiteness condition} we are assuming for all double complexes.

% \begin{lemma}\label{lemma:tot-is-exact}
%   The \hyperref[def:functor]{functor} \(\operatorname{tot}: \Complex[2]{\category{C}} \to \Complex{\category{C}}\) is \hyperref[def:exact-functor]{exact}.
% \end{lemma}
% \begin{proof}
%   Since \hyperref[ex:product-direct-sum-exact]{direct products are exact} and \(\operatorname{tot}(X^{\bullet, \bullet})^{n} = \bigoplus_{p+q=n}X^{p, q}\), the result follows.
% \end{proof}

% Recall that we \hyperref[rem:functor-double-complex-to-complex]{have defined functors \(F_{I}, F_{II}: \Complex[2]{\category{C}} \to \Complex{\Complex{\category{C}}}\)}.
% Since \(\category{C}\) is \hyperref[def:abelian-category]{category} we can combine these with the \hyperref[def:truncation-functor]{truncation functors}.
% For example we define \(\tau_{I}^{\leq n} \coloneq F_{I}^{-1} \circ \tau^{\leq n} \circ F_{I}\), i.e:
% \begin{equation*}
%   \tau_{I}^{\leq n}(X^{\bullet, \bullet}) \coloneq \cdots \to F_{I}(X^{\bullet, \bullet})^{n} \to \ker \d_{F_{I}(X^{\bullet, \bullet})}^{n} \to 0
% \end{equation*}
% and similary for \(\widetilde{\tau}_{I}^{\leq n}, \tau_{I}^{\geq n}, \widetilde{\tau}_{I}^{\geq n}\).

% For \(n \in \mathbb{Z}\) we define the \hyperref[def:complex]{complex}
% \begin{equation*}
%   \operatorname{H}_{I}^{n}(X^{\bullet, \bullet}) \coloneq \operatorname{H}^{n}(F_{I}(X^{\bullet, \bullet}))
% \end{equation*}
% and the \hyperref[def:double-complex]{double complex}
% \begin{equation*}
%   \operatorname{H}_{I}(X^{\bullet, \bullet}) \coloneq \cdots \to \operatorname{H}^{n}(F_{I}(X^{\bullet, \bullet})) \xrightarrow{0} \operatorname{H}^{n+1}(F_{I}(X^{\bullet, \bullet})) \to \cdots 
% \end{equation*}

% We get analogous definitions using \(F_{II}\) instead of \(F_{I}\).

% \begin{remark}
%   By \autoref{rem:truncation-functor-sequences} and \autoref{lemma:tot-is-exact} we get that the natural morphism
%   \begin{equation*}
%     \operatorname{tot}(\tau_{I}^{\leq n}(X^{\bullet, \bullet})) \to \operatorname{tot}(\widetilde{\tau}_{I}^{\leq n}(X^{\bullet, \bullet}))
%   \end{equation*}
%   is an \hyperref[def:categorical-isomorphism]{isomorphism} for all \(n \in \mathbb{N}\).

%   Furthermore we get the \hyperref[def:exact-sequence]{exact sequence} in \(\Complex{\category{C}}\)
%   \begin{center}
%     \begin{tikzcd}
%       0 \arrow[r] & \operatorname{tot}(\widetilde{\tau}_{I}^{\leq n-1}(X^{\bullet, \bullet})) \arrow[r] & \operatorname{tot}(\tau_{I}^{\leq n}(X^{\bullet, \bullet})) \arrow[r] & \operatorname{H}_{I}^{n}(X^{\bullet, \bullet})[-n] \arrow[r] & 0
%     \end{tikzcd}
%   \end{center}
%   which is \hyperref[def:functor]{functorial} with respect to \(X^{\bullet, \bullet}\).
% \end{remark}


% \begin{theorem}\label{thm:tot-is-qis}
%   Let \(f^{\bullet, \bullet}: X^{\bullet, \bullet} \to Y^{\bullet, \bullet}\) be a \hyperref[def:double-complex]{morphism of double complexes}.
%   If \(f^{\bullet, \bullet}\) induces an \hyperref[def:categorical-isomorphism]{isomorphism} \(\operatorname{H}_{II}\operatorname{H}_{I}(X^{\bullet, \bullet}) \cong \operatorname{H}_{II}\operatorname{H}_{I}(Y^{\bullet, \bullet})\) then \(\operatorname{tot}f: \operatorname{tot}(X^{\bullet, \bullet}) \to \operatorname{tot}(Y^{\bullet, \bullet})\) is a \hyperref[def:quasi-isomorphism]{quasi-isomorphism}.
% \end{theorem}
% \begin{proof}
%   Note that the hypothesis is equivalent to \(\operatorname{H}_{I}^{n}f: \operatorname{H}_{I}^{n}(X^{\bullet, \bullet}) \to \operatorname{H}_{I}^{n}(Y^{\bullet, \bullet})\) being a \hyperref[def:quasi-isomorphism]{quasi-isomorphism} for all \(n \in \mathbb{N}\).

%   Since by \autoref{rem:truncation-functor-sequences} we have
%   \begin{equation*}
%     \operatorname{H}_{I}^{n}(\tau_{I}^{\geq p}(X^{\bullet, \bullet})) \cong \begin{cases} \operatorname{H}_{I}^{n}(X^{\bullet, \bullet}) & n \geq p \\ 0 & n < p \end{cases}
%   \end{equation*}
%   for all \(p \in \mathbb{Z}\) we get the \hyperref[def:categorical-isomorphism]{isomorphism}
%   \begin{equation*}
%     \operatorname{H}_{II}\operatorname{H}_{I}(\tau_{I}^{\geq p}(X^{\bullet, \bullet})) \isoarrow \operatorname{H}_{II}\operatorname{H}_{I}(\tau_{I}^{\geq p}(Y^{\bullet, \bullet}))
%   \end{equation*}

%   For a fixed \(n \in \mathbb{Z}\) we have, for \(p \ll 0\)
%   \begin{equation*}
%     \operatorname{H}^{n}(\operatorname{tot}(X^{\bullet, \bullet})) \cong \operatorname{H}^{n}(\operatorname{tot}(\tau_{I}^{\geq p}(X^{\bullet, \bullet})))
%   \end{equation*}
%   and analogously for \(Y^{\bullet, \bullet}\).
%   Therefore by replacing \(X^{\bullet, \bullet}\) and \(Y^{\bullet, \bullet}\) by \(\tau_{I}^{\geq p}(X^{\bullet, \bullet})\) and \(\tau_{I}^{\geq p}(Y^{\bullet, \bullet})\) we may assume that \(X_{I}^{n} = Y_{I}^{n} = 0\) for \(n \ll 0\).

%   Consider now the commutative diagram
%   \begin{center}
%     \begin{tikzcd}
%       0 \arrow[r] & \operatorname{tot}(\widetilde{\tau}_{I}^{\leq n-1}(X^{\bullet, \bullet})) \arrow[r] \arrow[d, "{\operatorname{tot}(\widetilde{\tau}_{I}^{\leq n-1}f)}"'] & \operatorname{tot}(\tau_{I}^{\leq n}(X^{\bullet, \bullet})) \arrow[r] \arrow[d, "{\operatorname{tot}(\tau_{I}^{\leq n}f)}"] & \operatorname{H}_{I}^{n}(X^{\bullet, \bullet})[-n] \arrow[r] \arrow[d, "{(\operatorname{H}_{I}^{n}f)[-n]}"] & 0 \\
%       0 \arrow[r] & \operatorname{tot}(\widetilde{\tau}_{I}^{\leq n-1}(Y^{\bullet, \bullet})) \arrow[r] & \operatorname{tot}(\tau_{I}^{\leq n}(Y^{\bullet, \bullet})) \arrow[r] & \operatorname{H}_{I}^{n}(Y^{\bullet, \bullet})[-n] \arrow[r] & 0
%     \end{tikzcd}
%   \end{center}
%   Note that by hypothesis the right vertical arrow is a \hyperref[def:quasi-isomorphism]{quasi-isomorphism} and for \(n \ll 0\) the left arrow is also a \hyperref[def:quasi-isomorphism]{quasi-isomorphism}.
%   Therefore by induction and using \(\operatorname{tot}(\tau_{I}^{\leq n}(X^{\bullet, \bullet})) \cong \operatorname{tot}(\widetilde{\tau}_{I}^{\leq n}(X^{\bullet, \bullet}))\) we get that all vertical arrows are \hyperref[def:quasi-isomorphism]{quasi-isomorphisms}.

%   Finally the results follows since for a fixed \(n \in \mathbb{Z}\) and \(p \gg 0\) we have
%   \begin{equation*}
%     \operatorname{H}^{n}(\operatorname{tot}(X^{\bullet, \bullet})) \cong \operatorname{H}^{n}(\operatorname{tot}(\tau_{I}^{\leq p}(X^{\bullet, \bullet})))
%   \end{equation*}
% \end{proof}

% \begin{corollary}
%   Let \(X^{\bullet, \bullet}\) be a \hyperref[def:double-complex]{double complex} in \(\category{C}\).
%   If \(\operatorname{H}_{I}(X^{\bullet, \bullet}) = 0\) then \(\operatorname{tot}(X^{\bullet, \bullet})\) is \hyperref[def:quasi-isomorphism]{quasi-isomorphic to \(0\)}.
% \end{corollary}
% \begin{proof}
%   Take \(Y^{\bullet, \bullet} = 0\) in \autoref{thm:tot-is-qis}.
% \end{proof}

% \begin{corollary}\label{cor:tot-qis-to-zero}
%   Let \(X^{\bullet, \bullet}\) be a \hyperref[def:double-complex]{double complex} in \(\category{C}\).
%   If all rows \(X^{j, \bullet}\) are \hyperref[def:exact-sequence]{exact} for \(j \neq n\) then \(\operatorname{tot}(X^{\bullet, \bullet})\) is \hyperref[def:quasi-isomorphism]{quasi-isomorphic} to \(X^{n, \bullet}[-n]\).
% \end{corollary}
% \begin{proof}
%   Apply \autoref{thm:tot-is-qis} first to \(\sigma_{I}^{\geq n}(X^{\bullet, \bullet}) \to X^{\bullet, \bullet}\) and then to \(\sigma_{I}^{\geq n}(X^{\bullet, \bullet}) \to \sigma_{I}^{\leq n}\sigma_{I}^{\geq n}(X^{\bullet, \bullet}) \cong X^{n, \bullet}[-n]\).
% \end{proof}

% \begin{corollary}\label{cor:rows-cols-zero-and-exact}
%   Let \(X^{\bullet, \bullet}\) be a \hyperref[def:double-complex]{double complex} in \(\category{C}\).
%   If all rows \(X^{j, \bullet}\) and columns \(X^{\bullet, j}\) are zero for \(j < 0\) and \hyperref[def:exact-sequence]{exact} for \(j > 0\) then \(\operatorname{H}^{p}(X^{0, \bullet}) \cong \operatorname{H}^{p}(X^{\bullet, 0})\) for all \(p \in \mathbb{Z}\).
% \end{corollary}
% \begin{proof}
%   By hypothesis \(X^{0, \bullet}\) and \(X^{\bullet, 0}\) are \hyperref[def:quasi-isomorphism]{quasi-isomorphic} to 0.
% \end{proof}



\end{document}
